\PassOptionsToPackage{dvipsnames}{xcolor}
%\documentclass[12pt]{beamer}
%\documentclass[10pt]{beamer}
\documentclass[10pt, handout]{beamer}

\usepackage{mathtools} %% previously \usepackage{amsmath}
\usepackage{amsfonts}
\usepackage{amssymb} % http://www.ctan.org/tex-archive/macros/latex/required/psnfss/psnfss2e.pdf
%\usepackage[official]{eurosym}
\usepackage{helvet}
\usepackage{natbib}
\usepackage{lscape, longtable, array, booktabs, multirow, tabularx, graphics, epstopdf}
\usepackage{setspace}
\usepackage{listings}
\usepackage {tikz}

\usefonttheme[onlymath]{serif}

\usepackage{accents}
\newcommand*{\dt}[1]{%
   \accentset{\mbox{\LARGE $.$}}{#1}}

%\usepackage[T1]{fontenc}

\usecolortheme{dolphin}
%\useoutertheme{infolines}
\definecolor{dollarbill}{rgb}{0.52, 0.73, 0.4}
%\setbeamercovered{transparent}% Dim out "inactive" elements
\setbeamercovered{invisible}% Default
\beamertemplatenavigationsymbolsempty
\setbeamerfont{page number in head/foot}{size=\footnotesize}
\setbeamertemplate{footline}[frame number]


\usepackage{hyperref} 
\hypersetup{colorlinks=true}
%\hypersetup{linkcolor=RoyalBlue, urlcolor=ForestGreen, citecolor=Magenta}
\hypersetup{allcolors=RoyalBlue} %% 

%\renewcommand{\thefootnote}{\fnsymbol{footnote}}
\renewcommand{\thefootnote}{\arabic{footnote}}

\DeclareMathOperator*{\argmax}{argmax} % thin space, limits underneath in displays


\title[Linearization]{Linearization}
\vskip 0.1in

\subtitle{Mathematical Methods for Economics (771)}
%\author[H Hollander]{Hylton Hollander}

\institute[Stellenbosch University]{\large{
Stellenbosch University}}


%\date[April, 2020]{April, 2020}

\date{}

\begin{document}


\frame{\titlepage}

\frame{\frametitle{Contents:} \tableofcontents}


\section{Readings}

\frame{
    \frametitle{Readings} \vskip 0.1in


\begin{itemize}
\item Dejong, D. N. and C. Dave, 2007. Structural Macroeconometrics.
Princeton University Press. Chapter 2: Approximating and solving DSGE models;
\item A Toolkit for analyzing nonlinear dynamic stochastic
models easily (by Uhlig Harald).
\end{itemize}

}



\section{Introduction}

\begin{frame}

\frametitle{Introduction} \vskip 0.1in

A DSGE model consists of a system of nonlinear
difference equations. It typically includes three components (DeJong
and Dave (2007:11)):

\begin{enumerate}
\item the \textcolor{blue}{decision problem} -- a characterization of the environment in which decision makers
reside;
\item a set of \textcolor{blue}{decision rules} that dictate their behaviour (e.g., FOCs);
\item a characterization of the \textcolor{blue}{uncertainty} they face in making decisions (e.g., stochastic shocks, risk aversion \& information asymmetry).
\end{enumerate}

\end{frame}

\begin{frame}

\frametitle{Introduction} \vskip 0.1in

There generally exists no closed-form solution for such (non-linear) problems.
\begin{itemize}
\item[$\therefore$] we must result to numerical and/or approximation techniques $\rightarrow$ construct a linear approximation of the model.
\end{itemize}\vskip 0.1in

{{There are varied approaches to approximate a nonlinear system
linearly {(\footnotesize{see, e.g., Judd (1998).)}}, such as
\begin{itemize}
\item local approximation (eg. Taylor series and Pad{\'{e}}
expansions);
\item interpolation;
\item regression.
\end{itemize}}}\vskip 0.1in

In this session, we study the log-linearization method only.
\begin{itemize}
  \item the standard method;
  \item a simpler method (Uhlig's method);
\end{itemize}

\end{frame}


\begin{frame}

\frametitle{Introduction} \vskip 0.1in

We first:\vskip 0.05in
\begin{enumerate}
\item[(1)] take natural logs of the system of non-linear difference equations
\end{enumerate}\vskip 0.1in

We then:\vskip 0.05in
\begin{enumerate}
\item[(2)] linearize the logged difference equations about a particular point (usually a steady state). And simplify until we have a system of linear difference equations:\vskip 0.05in
    \begin{enumerate}
    \item[$\triangleright$] variables are expressed as (percentage) deviations about their steady state values;\vskip 0.05in
    \item[$\triangleright$] directly amenable to empirical implementation
    \end{enumerate}
\end{enumerate}

\end{frame}

\begin{frame}
\frametitle{Introduction} \vskip 0.1in

\textcolor{blue}{Why this method?} \vskip0.1in

\begin{itemize}
\item Linearization is nice because we know how to work with linear difference equations;
\item Log-difference terms are nice because they provide natural interpretations of the units (i.e., everything is in percentage terms)
\end{itemize}

\vskip0.2in

{\footnotesize{Since the principle of both methods that we focus on here is to use
a Taylor-series expansion around the steady state to replace all
equations by approximations, let's first review the theorem of
Taylor series approx.}}

\end{frame}

\section{Taylor Series Approximation}

\begin{frame}

\frametitle{Taylor Series Approximation}


\textcolor{blue}{Taylor's theorem for a univariate function $f(x)$:}
\begin{eqnarray}
f(x) &=& f(x^*) + \frac{f'(x^*)}{1!}(x-x^*)+ \frac{f''(x^*)}{2!}(x-x^*)^2  \nonumber\\ &+& \frac{f^{(3)}(x^*)}{3!}(x-x^*)^3 +  \ldots  + \frac{f^{n}(x^*)}{n!}(x-x^*)^n +
R_{n+1}(x)\label{EqB.1}%\\
%R_{n+1}(x) &=& \int_{x^*}^{x} \frac{(x-t)^n}{n!} f^{n+1} (t)dt
%\label{EqB.2}
\end{eqnarray}
{\footnotesize{where we refer to $R_{n+1}(x)$ as the residual term.}}\vskip 0.2in

{\footnotesize{Taylor's theorem says that one can use derivative information at a
single point to construct a polynomial approximation of a function
at a point ($x^*$)---usually a steady state. It is the theoretical basis for all local
approximations.}}

\end{frame}

%\begin{frame}

%\frametitle{Taylor Series Approximation} \vskip 0.1in

%Following Dejong and Dave (2007), consider we have a n-equation
%system of nonlinear difference equations:

%\begin{equation}\label{EqB.3}
%\psi (z_{t+1},z_t) = 0
%\end{equation} \pause

%By applying the first-order Taylor Series approximation, we can
%linearize \eqref{EqB.3} around its steady state $\psi(\bar{z})$,

%\begin{equation}\label{EqB.4}
%0 \approx \psi(z^{ss}) + \frac{\partial \psi}{\partial
%z_t}(z^{ss})(z_t - z^{ss}) + \frac{\partial \psi}{\partial
%z_{t+1}}(z^{ss})(z_{t+1} - z^{ss})
%\end{equation}

%where $\frac{\partial \psi}{\partial z_t}(z^{ss})$ denotes the
%Jacobian of $\psi (z_{t+1},z_t)$ w.r.t. $z_t$ at $z^{ss}$.\vskip
%0.1in

%\pause

%The solution of \eqref{EqB.3} is given by,

%\begin{equation}\label{EqB.5}
%A (x_{t+1}) = B (x_t)
%\end{equation}

%where $A=\frac{\partial \psi}{\partial z_{t+1}}(z^{ss})$,
%$B=-\frac{\partial \psi}{\partial z_t}(z^{ss})$, $x_t=z_{t+1}-z_t$
%\end{frame}

\begin{frame}

\frametitle{Taylor Series Approximation} \vskip 0.1in

\textcolor{blue}{We go one step further:}\vskip0.05in

For a function that is sufficiently smooth, the higher order derivatives will be small, and the function can be well approximated (about $x^*$) linearly as:
\begin{equation}
f(x) \approx f(x^*) + f'(x^*)(x - x^*)
\end{equation}

\vskip0.1in
\pause

\textcolor{blue}{The theorem also applies to multivariate functions:}\vskip0.05in

\noindent For example, the first order approximation of $f(x,y)$ about the point $(x^*,y^*)$ is:
\begin{equation}
f(x,y) \approx f(x^*,y^*) + f_{x}(x^*,y^*)(x - x^*) + f_{y}(x^*,y^*)(y - y^*)
\end{equation}

\end{frame}


\section{The Standard Method}

\begin{frame}

\frametitle{The Standard Method} \vskip 0.1in

{\footnotesize{The log-linearization method is based on the assumption that if the
system is not too far from the steady state, the linear version
result closely approximates the original system.}}\vskip 0.1in

Consider a model that consists of a set of nonlinear equations of
the following general form,

\begin{equation}\label{EqB.6}
F(x_{t})=\frac{G(x_t)}{H(x_t)}
\end{equation}

Taking the logarithms,

\begin{equation}\label{EqB.7}
ln(F(x_{t}))=ln(G(x_t))- ln(H(x_t))
\end{equation}

Again, by applying the first-order Taylor Series approximation, we
can linearize \eqref{EqB.7} around its steady state $x^{ss}$.


\end{frame}

\begin{frame}

\frametitle{The Standard Method} \vskip 0.1in

The log linearized \eqref{EqB.7},

{\scriptsize{
\begin{eqnarray}\label{EqB.8}
ln(F(x^{ss})) + \frac{F'(x^{ss})}{F(x^{ss})}(x_t-x^{ss}) &=&
[ln(G(x^{ss}) + \frac{G'(x^{ss})}{G(x^{ss})}(x_t-x^{ss})] -
[ln(H(x^{ss})) \nonumber\\
&& + \frac{H'(x^{ss})}{H(x^{ss})}(x_t-x^{ss})]
\end{eqnarray}}}

\pause

To simplify the model, we can eliminate the steady state of the
model from \eqref{EqB.8},

\begin{equation}\label{EqB.9}
\frac{F'(x^{ss})}{F(x^{ss})}(x_t-x^{ss}) =
\frac{G'(x^{ss})}{G(x^{ss})}(x_t-x^{ss}) -
\frac{H'(x^{ss})}{H(x^{ss})}(x_t-x^{ss})
\end{equation}

The model is now linear in $x_t$.
\end{frame}


\begin{frame}

\frametitle{The Standard Method} \vskip 0.1in

The above procedure applies equally well in multivariate cases. To summarize, the 3-step procedure for log-linearizing is:\vskip0.1in

\begin{enumerate}
\item Take logs
\item Do a first order Taylor series expansion about a point (the steady state)
\item Simplify so that everything is expressed in {$\%$} deviations from steady state
\end{enumerate}

\end{frame}


\section{Uhlig's Method}

\begin{frame}

\frametitle{Uhlig's Method} \vskip 0.1in

Uhlig's method does not require explicit differentiation and the
linearized model is expressed in terms of log differences of the
variables.\vskip 0.1in

\pause

Suppose $X_t$ be the vector of variables, $X^{ss}$ their steady
state, and $x_t$ the vector of log-deviations:

\begin{equation}\label{EqB.10}
x_t=ln{X_{t}}-ln{X^{ss}}
\end{equation}

That is, $x_t$ denotes the percentage deviation from their steady
state levels. Alternatively\footnote{Given \eqref{EqB.10},
$x_t=ln\frac{X_t}{X^{ss}}$, $e^{x_t}=\frac{X_t}{X^{ss}}$,
$\therefore X_t = X^{ss}e^{x_t}$; $
e^{x_t}=\sum_{n=1}^{\infty}\frac{(1+x_t)^n}{n!} \approx 1+x_t$.}

\begin{equation}\label{EqB.11}
X_t=X^{ss}e^{x_t}\approx X^{ss}(1+x_t)
\end{equation}

\vfill

{\footnotesize 1. Given \eqref{EqB.10} }


\end{frame}

\begin{frame}

\frametitle{Uhlig's Method} \vskip 0.1in

\textcolor{blue}{Four steps (Uhlig, 1995:5):}

\begin{enumerate}
\item multiply out everything in each equation;
\item replace a variable $X_t$ with $X_t=X^{ss}e^{x_t}$;
\item take logarithms, where both sides of an equation only involve
products, or use the following three building blocks\footnote{Eg.
$e^{x_t}\approx 1+x_{t}$; $aX_{t}=aX^{ss}e^{x_{t}}$;
$(X_t+a)Y_t=X^{ss}Y^{ss}e^{x_t}+X^{ss}Y^{ss}e^{y_t}+aY^{ss}y_t$}:
\begin{eqnarray}
e^{x_{t}+ay_{t}} &\approx& 1+x_t+ay_t\\
x_t y_t &\approx& 0\\
E_{t}[ae^{x_{t+1}}] &\approx& E_{t}[ax_{t+1}]
\end{eqnarray}
\item drop constants in each equation in the end.
\end{enumerate}

\vfill

{\footnotesize {2. E.g.,
$e^{x_t}\approx 1+x_{t}$; $aX_{t}=aX^{ss}e^{x_{t}}$;
$(X_t+a)Y_t=X^{ss}Y^{ss}e^{x_t}+X^{ss}Y^{ss}e^{y_t}+aY^{ss}y_t$}}


\end{frame}


\section{Examples}

\begin{frame}

\frametitle{Examples} 

Log-linearize the following using both the Taylor method and Uhlig's method, and compare your answers: \vskip0.2in

(1) The resource constraint $$Y_t = C_t + I_t ~,$$ indicating that output can be either consumed or invested.\vskip0.1in

(2) The Cobb-Douglas production function $$Y_t=a_{t}K^{\alpha}_{t}N^{1-\alpha}_t ~,$$ where output is produced by use of capital and labour and is subject to a productivity shock.

\end{frame}



\begin{frame}

\frametitle{Hansen's real business cycle model (Uhlig, pp. 5-7)} 

With the tools you've learned from dynamic programming, you can now simulate a dynamic stochastic general equilibrium model. \vskip0.1in 

We will work use Hansen's real business cycle model (Uhlig, pp. 5-7). \vskip0.1in 

See the tutorial on recursive methods and linearization on SUNLearn. \vskip0.1in

For the final question you will be given step-by-step instructions to simulate the model and obtain impulse response functions (of the variables) in response to a 1 percentage point increase in total factor productivity (see figure in the next slide). \vskip0.1in

Both Dynare and Octave are free software. Octave is basically a free version of matlab.

\end{frame}


\begin{frame}
\frametitle{Impulse response functions}
\begin{figure}[hpb]\centering \vskip-1in
\includegraphics[width=0.75\textwidth]{basicrbc_fig.pdf}
%\caption{IRFs to a 1 percentage point increase in total factor productivity}
\end{figure}
\end{frame}


%\begin{frame}
%\begin{center}
%\Huge{THE END}
%\end{center}
%\end{frame}
\end{document}
